\section{Observation spaces and substitution-stable candidates}
\label{obs:section}

We now construct the candidate interpretation for the remaining primary
component.  Throughout this section, $S$ is the current component of the
primary graph and $U$ is its predecessor-closed processed part.  We retain
the hypothesis of weak normalization for every legal expression.  We use
the preceding profile and candidate lemmas, including
$G_T=\operatorname{Ty}(T)\cap\mathrm{SN}_\beta$ when $\pi(T)\in U$.
The new input will be the exclusion of a particular finite graph
configuration, stated below and proved in the next part of the argument.
No strong-normalization hypothesis is imposed on substitutions in this
section.

The interpretation of a neutral type must account for substitutions
that replace its head by a product.  We therefore give variables parameters
that record candidates, and interpret the syntax of a substituted term
to determine which candidate to read.  At a function type
$T=\Pi x:D.E$, such a parameter records a base observation.  The
construction supplies spaces only at selected types; whenever the
instantiated result $(E[q/x])^\#$ has a space, a normal argument
$q:D$, together with a parameter for $q$ when its type requires one,
selects a next value in that result space.  The function's space can
therefore depend on the spaces for both its arguments and its results.  These dependencies
can cycle, so ordinary recursion on the displayed type does not suffice.

The graph below identifies which spaces are needed and separates two
ways to construct them.  Some admit a description by finite observation
paths or a decreasing count of type nodes.  For the others we build two
spaces together: unrestricted tables on one side and tables with a
uniform finite bound on their payload dependence on the other.  After
constructing the spaces, we prove that every finite typed syntax tree
can be evaluated in them.  Substitution and reduction require an
additional condition, called goodness, which holds for the normal
syntax used to define our candidates.  That final transport result will
supply the specialization property of
the candidate interface.

\subsection{Which values require observations?}
\label{obs:graphs}

There are two different profiles to keep apart.  A neutral type sorted
at $s$ has expression profile $\{s\}$; a value interpreting that type
is indexed by its expected type $s$, whose profile is
$\operatorname{Ax}(s)$.  These axiom profiles are the starting points
for deciding which parameter spaces are needed.  The remaining graph
definitions follow the dependencies of those spaces toward
domains that are tested by products within $S$.

Let $H$ be the closure under primary positive paths of
\[
 \{\operatorname{Ax}(s):\{s\}\in S,
                         \ \operatorname{Ax}(s)\ne\varnothing\}.
\]
The \emph{secondary graph} has vertex set $H$.  It retains all primary
positive edges between these vertices.  It retains the negative edge
$I\longrightarrow K$ belonging to a profile triple $(I,J,K)$ precisely
when $I,J,K\in H$.  A vertex $I\in H$ is \emph{direct} if some profile
triple $(I,J,K)$ has $J,K\in S$; here $J,K$ need not belong to $H$.
A vertex is \emph{active} if it has a secondary path to a direct vertex.
Only the induced graph on active vertices will be used below.

An active vertex is \emph{plain} if paths from it to direct vertices
occur with both parities.  Otherwise it is \emph{free}; its sign is $+$
if those paths are even and $-$ if they are odd.  A path of length zero
is even, so a free direct vertex has sign $+$.

This last observation will control the later uniform bounds.  A product
test whose product and instantiated codomain profiles remain in $S$
has its open codomain profile in $S$ as well, by their positive ascent.
A domain profile that is active is therefore direct.  Such a test may
quantify over parameters at plain or free plus profiles, but never at
free minus profiles.
This fact will let one finite dependence bound survive all the parameter
choices in a product test.

\begin{lemma}[Signs and dependencies]\label{obs:graph-lemma}
An active predecessor of a plain vertex is plain.  Along an edge between
free vertices, the sign is preserved on a positive edge and reversed on
a negative edge.

Let $T=\Pi x:D.E$ be a normalized actual type, put
$I=\pi(D)$, $J=\pi(E)$, and $K=\pi(T)$, and let
$q\in\operatorname{Ty}(D)$.  Put $T_q=(E[q/x])^\#$ and $L=\pi(T_q)$.
If $K,L$ are active, then $J$ is active, and $I$ is active whenever
$I\in H$.  If instead $K,L\in S$, then $J\in S$ and every active $I$
is direct.
\end{lemma}
\begin{proof}
Prepending the edge in question to paths to a direct vertex proves the
sign statements.  Profile growth and the product edge give the positive
path
\[
                         L\longrightarrow J\longrightarrow K,
\]
where the first edge is omitted when $L=J$.  Since $L\in H$ and $H$ is
positive-path closed, $J\in H$.  Its path to the active vertex $K$ makes
$J$ active.  If $I\in H$, the negative edge of $(I,J,K)$ is retained,
so $I$ is active too.  When $K,L\in S$, the displayed path cannot leave
their primary strongly connected component.  Thus $J\in S$, and the
definition of direct applies.
\end{proof}

The \emph{layer} of a type $T$ is $\pi(T)$.  A type is active, plain, or
free of a given sign when its layer is so; an empty layer is inactive.
These terms concern the type of a value, not the profile of the value
itself.

Here is the promised graph hypothesis.
\begin{equation}\label{obs:forbidden}
\begin{gathered}
\text{There is no plain active secondary strongly connected component $C$}\\[-2pt]
\text{with an internal negative edge, a profile triple $(I,J,K)$ with $J,K\in C$,}\\[-2pt]
\text{and a sort $s$ such that $\{s\}\in C$, $\operatorname{Ax}(s)\ne\varnothing$, and}\\[-2pt]
                  \operatorname{Ax}(s)\longrightarrow I
                  \text{ by a primary positive path.}
\end{gathered}
\end{equation}
All results in this section are conditional on
\eqref{obs:forbidden}.  The later typed contradiction will derive this
condition from the same system-wide weak-normalization hypothesis.

\subsection{Actual frames and plain observation spaces}
\label{obs:plain-spaces}

Let $\mathcal B$ be the set of assignments choosing one candidate at
every actual type.  Its distinguished member $\mathbf t$ chooses every
top candidate.  All types and expressions here are syntax over finite
prefixes of $\Omega$, modulo $\alpha$-equivalence, so their collections
and $\mathcal B$ are sets.

We construct a nonempty set $\mathcal K_T$ of parameters for every
active actual type $T$.  For a plain type we require the following
precise prefix description.  If $T=\Pi x:D.E$, a \emph{legal frame at
$T$} consists of a normal term $q\in\operatorname{Ty}(D)$ for which
$T_q=(E[q/x])^\#$ is active, together with a parameter
$v\in\mathcal K_D$ if $D$ is active.  If $D$ is inactive there is one
absent-payload symbol in place of $v$.  Its target is $T_q$.  A
nonproduct type has no legal frames.  Write $\mathcal F_T$ for this
set of frames.  The desired description is the bijection
\begin{equation}\label{obs:plain-prefix}
 \mathcal K_T\ \simeq\
 \mathcal B\times
       \prod_{a\in\mathcal F_T}\mathcal K_{\operatorname{target}(a)}.
\end{equation}
We call its first coordinate the base and its $a$-coordinate the next
value.  By Lemma~\ref{obs:graph-lemma}, all active spaces occurring on
the right are plain when $T$ is plain, and their secondary components
precede or equal that of $T$.

\begin{lemma}[Construction of plain spaces]\label{obs:plain-construction}
Under \eqref{obs:forbidden}, the plain spaces admit
\eqref{obs:plain-prefix}.  They also admit distinguished values whose
bases are $\mathbf t$ and whose next values are distinguished values.
\end{lemma}
\begin{proof}
Process the finitely many plain secondary components in source order.
Suppose first that a component $C$ has no internal negative edge.
Every active domain used by a legal frame at a type of layer in $C$
then lies in an earlier component: otherwise the frame's triple, whose
open codomain is active, would contribute an internal negative edge.
Thus all frame alphabets can be specified without using a space whose
layer belongs to $C$.

For $\pi(T)\in C$, consider finite paths of legal frames beginning at
$T$.  A continuing path has all its target layers in $C$, and includes
the empty path.  An exiting path has all but its final target layer in
$C$ and its final target in an earlier component.  Define $\mathcal K_T$
to be the set of tables assigning a member of $\mathcal B$ to each
continuing path and a member of the final target space to each exiting
path.  Split a table according to the empty path and the first frame.
For a frame staying in $C$, its suffix table is exactly a table of the
same kind starting at its target.  For an exiting frame its output is
already a value of the target space.  This gives the bijection
\eqref{obs:plain-prefix} in both directions.  The constant top outputs
and previously chosen exit defaults define the distinguished table.

Now let $C$ have an internal negative edge.  For every normalized type
$T$ of layer in $C$, define a positive integer $p_C(T)$ as follows.
Count its root.  At a product with child layers $I,J$, continue into
the codomain when $J\in C$, and into the domain when $I\in C$ and
$J$ is active.  Apply the same rule recursively at counted children,
using the binder context in the codomain.  This counts a subset of a
finite syntax tree.  Profile invariance makes the count independent of
the prefix and unchanged by valid thinning or context conversion.

For a legal frame at $T=\Pi x:D.E$, write
$I=\pi(D)$, $J=\pi(E)$, $K=\pi(T)$, and $L=\pi(T_q)$.
A same-$C$ domain has smaller count.  We claim
that its same-$C$ target also satisfies
\begin{equation}\label{obs:potential-drop}
               p_C(T_q)\le p_C(E)<p_C(T).
\end{equation}
Here $J\in C$ follows from the positive ascent in
Lemma~\ref{obs:graph-lemma}; hence $p_C(E)$ is defined.

To prove the first inequality, thin $T$ and $q$ to a common prefix,
choosing the displayed binder fresh.  Associate to every counted node
of $(E[q/x])^\#$ its position string.  We show that the same position
is a counted node of $E$.  The induction along a position maintains
that the corresponding pre- and post-substitution parent profiles both
lie in $C$; at the root these are $J$ and $L$.  At a pre-substitution
product the post normal form is a product with componentwise substituted
normal forms.  Write $I_0,J_0,K_0$ for its pre domain, codomain, and
parent profiles, and $I_1,J_1,K_1$ for the post profiles.  Profile growth
gives $I_1\supseteq I_0$ and $J_1\supseteq J_0$, also under binders
after context conversion.

If the post codomain is counted, then $J_1\in C$, and the positive
path $J_1\longrightarrow J_0\longrightarrow K_0$ has both endpoints
in $C$.  Hence $J_0\in C$, so the pre codomain is counted too.  If
the post domain is counted, then $I_1\in C$ and $J_1$ is active.
The inclusion paths put $I_0,J_0$ in $H$.  Since $K_0\in C\subseteq H$,
the pre triple has a retained negative edge $I_0\longrightarrow K_0$.
The path
\[
                  I_1\longrightarrow I_0\longrightarrow K_0
\]
again has endpoints in $C$, so $I_0\in C$.  Also the positive edge
$J_0\longrightarrow K_0$ makes $J_0$ active.  Thus the pre domain is
counted.  In either case the pre- and post-substitution profiles of the
selected child lie in $C$, which preserves the induction invariant for
the next position.  Inclusion edges in these paths are omitted when
the profiles are equal.

The association could fail only if a pre neutral leaf became a product.
Sorts are unchanged, and a neutral with any other head remains neutral
after substitution and normalization.  Its head must therefore be the
substituted variable $x$.  Let its singleton profile be $\{s\}$;
this lies in $C$.  The sort $s$ is sorted, so
$\operatorname{Ax}(s)\ne\varnothing$.  Trace the original $x$-headed
spine backward through its normalized successive typing products.
At each application the result layer has a primary positive path to
the function layer.  Starting with the layer
$\operatorname{Ax}(s)$ of its type $s$, concatenate these paths to
obtain a positive path to $\pi(D)=I$.  The outer product has
$J,K\in C$.  This is forbidden by \eqref{obs:forbidden}.  No leaf
therefore expands.  Position strings give the required injection of
counted nodes and prove \eqref{obs:potential-drop}.

Define \eqref{obs:plain-prefix} by induction on $p_C(T)$, simultaneously
for all types with the same count.  Every same-component space on its
right is already constructed by the strict inequalities just proved;
all other spaces come from earlier components.  The product of the
specified nonempty sets is nonempty, with the explicit element whose
base is $\mathbf t$ and whose next values are the already specified
defaults.  This completes the construction.
\end{proof}

\subsection{Free observation spaces}
\label{obs:free-spaces}

The construction in this subsection is purely set-theoretic.  The typing
conditions on frames will be imposed by evaluation; they do not restrict
the tables constructed here.  Fix the set $\mathcal Q$ of normal syntax
over $\Omega$, with bound names identified up to $\alpha$-conversion.
Let $\mathcal T_{\mathrm{pl}}$ be the set of actual plain active types,
and suppose that the nonempty spaces $\mathcal K_D$ have already been
constructed for $D\in\mathcal T_{\mathrm{pl}}$.  The fixed set of tagged
plain values is
\begin{equation}
  \mathcal V_{\mathrm{pl}}
   =\coprod_{D\in\mathcal T_{\mathrm{pl}}}
       \bigl(\{D\}\times\mathcal K_D\bigr).
  \label{obs:free-plain-values}
\end{equation}
Thus a member of $\mathcal V_{\mathrm{pl}}$ remembers both its type $D$ and its value
$v\in\mathcal K_D$.  Let $\mathcal B$ be the set of base observations:
a member chooses a candidate at every actual type.  Write
$\mathbf t\in\mathcal B$ for the choice of the top candidate everywhere.
For the present construction only the set $\mathcal B$ and its
distinguished member $\mathbf t$ are needed.  Put
\[
  \mathcal O=\{\star\}\sqcup\mathcal V_{\mathrm{pl}},
\]
where the union is disjoint.  In particular, both output sets
$\mathcal B$ and $\mathcal O$ have specified default elements, even if
there are no plain active types.

\paragraph{Frames, stacks, and tables.}
For any set $Y$ define the payload set
\[
  V(Y)
   =\{\mathrm{none}\}
      \sqcup\bigl(\{\mathrm{opp}\}\times Y\bigr)
      \sqcup\bigl(\{\mathrm{plain}\}\times\mathcal V_{\mathrm{pl}}\bigr).
\]
The three tags make these alternatives disjoint.  Define continue and
exit frames by
\[
  F_c(Y)=\mathcal Q\times V(Y)\times\{c\},
  \qquad
  F_e(Y)=\mathcal Q\times V(Y)\times\{e\}.
\]
Here $Y$ is the space of opposite-sign payloads.  The set of non-exiting
stacks and the set of exiting stacks are, respectively,
\begin{equation}
  S_c(Y)=\coprod_{n\geq 0} F_c(Y)^n,
  \qquad
  S_e(Y)=\coprod_{n\geq 0}
                 \bigl(F_c(Y)^n\times F_e(Y)\bigr).
  \label{obs:free-stacks}
\end{equation}
Thus $S_c(Y)$ contains the empty stack $\epsilon$, and an exit frame
occurs exactly once, at the end of a stack in $S_e(Y)$.  Set
\begin{equation}
  \operatorname{Tab}(Y)
      =\mathcal B^{S_c(Y)}\times\mathcal O^{S_e(Y)}.
  \label{obs:free-tab}
\end{equation}
A table $t=(t_c,t_e)$ gives a base observation on each non-exiting
stack and an exit observation on each exiting stack.  Its default
$d_Y$ is specified by
\[
  (d_Y)_c(\sigma)=\mathbf t,
  \qquad
  (d_Y)_e(\tau)=\star.
\]
Consequently $\operatorname{Tab}(Y)$ is nonempty for every $Y$,
including $Y=\emptyset$.

We seek two spaces whose tables take values of the opposite sign as
payloads.  Their construction must solve two different problems.  A plus
value should accept every finite stack of minus values.  A minus value
must be representable at a finite stage, even though it is tested on
arbitrarily many plus values and arbitrarily long stacks.  Choosing a
stage separately for each stack would not give that representation.

The construction below supplies a space $C_*$ of minus values and the
full space $P_*=\operatorname{Tab}(C_*)$ of plus tables.  Increasing
stages $C_r$ give restrictions
$\rho_r:P_*\to P_r=\operatorname{Tab}(C_r)$ of plus tables.
A table on plus payloads represents a minus value precisely when
one index $r$ works for every observation: replacing each plus payload
$p$ by $p'$ with $\rho_r(p)=\rho_r(p')$ leaves the table output
unchanged.  Finally, fixing any continue prefix must again give a value
of the same sign; on the minus side it must retain the same representing
stage.  These are the representation and slicing properties that typed
evaluation will use.  We first construct the stages and then prove this
characterization.

\paragraph{Relabeling and contravariance.}
A map $f:Y\to Z$ induces stack maps
$S_a(f):S_a(Y)\to S_a(Z)$ for $a\in\{c,e\}$.
These replace $(\mathrm{opp},y)$ by $(\mathrm{opp},f(y))$ in each
payload position.  Keys, flags, and the other payload alternatives
are unchanged.  Define
\begin{equation}
  \operatorname{Tab}(f):\operatorname{Tab}(Z)
                  \longrightarrow\operatorname{Tab}(Y),
  \qquad
  \bigl(\operatorname{Tab}(f)(t)\bigr)_a
                   =t_a\circ S_a(f).
  \label{obs:free-contravariance}
\end{equation}
Relabeling preserves identities and composition.  Hence
\[
  \operatorname{Tab}(\operatorname{id}_Y)
     =\operatorname{id}_{\operatorname{Tab}(Y)},
  \qquad
  \operatorname{Tab}(g\circ f)
     =\operatorname{Tab}(f)\circ\operatorname{Tab}(g).
\]
It also preserves the default tables:
$\operatorname{Tab}(f)(d_Z)=d_Y$.

We will use the following two elementary facts.
\begin{enumerate}
\item If $f$ is injective, then $S_c(f)$ and $S_e(f)$ are injective,
and $\operatorname{Tab}(f)$ is surjective.  Indeed, given a table on
$Y$, define a table on $Z$ by transporting its values to the images
of $S_c(f)$ and $S_e(f)$, and assigning $\mathbf t$ or $\star$ on the
respective complements.  Injectivity makes the transported values
unambiguous, and restriction recovers the given table.
\item If $f$ is surjective, then $S_c(f)$ and $S_e(f)$ are surjective,
and $\operatorname{Tab}(f)$ is injective.  To lift a stack, lift its
finitely many opposite payloads through $f$.  If two tables become
equal after precomposition with these surjective stack maps, they
agree on every stack and therefore are equal.
\end{enumerate}
In particular these statements apply when an opposite payload set is
empty; the stacks without opposite payloads remain available.

\paragraph{The alternating chain.}
Define, for $r=0,1,2,\ldots$,
\begin{equation}
  C_0=\emptyset,
  \qquad P_r=\operatorname{Tab}(C_r),
  \qquad C_{r+1}=\operatorname{Tab}(P_r).
  \label{obs:free-stages}
\end{equation}
Let $i_0:C_0\to C_1$ be the empty map.  Inductively, once
$i_r:C_r\to C_{r+1}$ is defined, put
\begin{equation}
  j_r=\operatorname{Tab}(i_r):P_{r+1}\longrightarrow P_r,
  \qquad
  i_{r+1}=\operatorname{Tab}(j_r):C_{r+1}\longrightarrow C_{r+2}.
  \label{obs:free-bonding}
\end{equation}
The preceding facts prove by induction that every $i_r$ is injective
and every $j_r$ is surjective.  Each $j_r$ is restriction to stacks
whose opposite payloads lie in the embedded copy of $C_r$; its
surjectivity has the explicit default-extension proof above.

For $r\leq s$, write $i_{r,s}:C_r\to C_s$ for the composite of the
injections, with $i_{r,r}$ the identity.  Let $C_*$ be the set colimit
of this chain.  Explicitly, take the disjoint union of the $C_r$ and
identify $(r,x)$ with $(s,y)$ precisely when, at some common stage
$m\geq r,s$,
\[
  i_{r,m}(x)=i_{s,m}(y).
\]
The canonical maps $e_r:C_r\to C_*$ are injective and satisfy
$e_{r+1}\circ i_r=e_r$.  Their images form an increasing union
equal to $C_*$.  We may therefore regard the stages as nested
subsets after these identifications.  Every finite family of elements
of $C_*$ lies in a common stage: choose a representing stage for
each element and then take their maximum.  Every element in fact
has a representative in some $C_{r+1}$, since $C_0$ is empty.

\paragraph{Final plus tables and finite-stage agreement.}
Put
\begin{equation}
  P_*=\operatorname{Tab}(C_*),
  \qquad
  \rho_r=\operatorname{Tab}(e_r):P_*\longrightarrow P_r.
  \label{obs:free-final-plus}
\end{equation}
Each $\rho_r$ is surjective, by default extension along the injection
$e_r$.  Contravariance and the identities for $e_r$ give
\begin{equation}
  \rho_r=j_r\circ\rho_{r+1}.
  \label{obs:free-restriction-compatibility}
\end{equation}
For $p,p'\in P_*$ write
\begin{equation}
  p\equiv_r p' \quad\Longleftrightarrow\quad
           \rho_r(p)=\rho_r(p').
  \label{obs:free-equivalence}
\end{equation}
Agreement at a larger stage implies agreement at every smaller
stage, by repeated use of
\eqref{obs:free-restriction-compatibility}.

Every stack in $S_a(C_*)$, for $a\in\{c,e\}$, belongs to the image
of $S_a(e_r)$ for some $r$: it has only finitely many opposite
payloads, which lie in a common stage.  The same holds for any
finite family of stacks, with one common $r$.  It follows that
\begin{equation}
  \bigl(p\equiv_r p'\text{ for every }r\bigr)
       \quad\Longleftrightarrow\quad p=p'.
  \label{obs:free-separation}
\end{equation}
For the forward implication, lift any specified final stack from a
common stage and evaluate the equality $\rho_r(p)=\rho_r(p')$ on
that lift.  The reverse implication is immediate.  No assertion
that a plus table itself belongs to a finite stage is needed.

\paragraph{Final minus behaviors.}
For each $r\geq0$, define
\begin{equation}
  \beta_r=\operatorname{Tab}(\rho_r):
        C_{r+1}=\operatorname{Tab}(P_r)
            \longrightarrow\operatorname{Tab}(P_*).
  \label{obs:free-behavior-stage}
\end{equation}
In coordinates, for $u\in C_{r+1}$ and $\sigma\in S_a(P_*)$,
\[
  \bigl(\beta_r(u)\bigr)_a(\sigma)
       =u_a\bigl(S_a(\rho_r)(\sigma)\bigr).
\]
Thus each plus payload of a final stack is first restricted to
$P_r$.  Since $\rho_r$ is surjective, $\beta_r$ is injective.
Moreover,
\begin{align*}
  \beta_{r+1}\circ i_{r+1}
    &=\operatorname{Tab}(\rho_{r+1})
                          \circ\operatorname{Tab}(j_r)\\
    &=\operatorname{Tab}(j_r\circ\rho_{r+1})
      =\operatorname{Tab}(\rho_r)=\beta_r.
\end{align*}
These maps consequently induce a well-defined map
\begin{equation}
  \beta:C_*\longrightarrow\operatorname{Tab}(P_*),
  \qquad
  \beta(e_{r+1}(u))=\beta_r(u).
  \label{obs:free-behavior}
\end{equation}
It is injective: represent any two proposed equal-behavior elements
at a common stage $C_{r+1}$ and use the injectivity of $\beta_r$.
We henceforth evaluate a minus value through this behavior map.

The defaults are compatible with all bonding maps.  In particular,
the default elements $d_{P_r}\in C_{r+1}$ represent one element of
$C_*$, whose behavior is $d_{P_*}$.  The plus default is
$d_{C_*}\in P_*$.  Both free spaces are therefore nonempty.

\paragraph{Exact characterization by a uniform stage bound.}
Fix $r$.  For $a\in\{c,e\}$ and $\sigma,\tau\in S_a(P_*)$ write
$\sigma\sim_r\tau$ when they have identical lengths, keys, flags,
and payload alternatives, their plain payloads are identical, and
each corresponding pair of opposite payloads $p,p'$ satisfies
$p\equiv_r p'$.  Equivalently,
\begin{equation}
  \sigma\sim_r\tau
    \quad\Longleftrightarrow\quad
  S_a(\rho_r)(\sigma)=S_a(\rho_r)(\tau).
  \label{obs:free-stack-equivalence}
\end{equation}
This equivalence follows directly from the disjoint payload tags
and the definition of stack relabeling.

For a final table $t\in\operatorname{Tab}(P_*)$, the following
conditions are equivalent:
\begin{enumerate}
\item $t$ belongs to the image of $\beta_r$.
\item For both $a=c$ and $a=e$, the equality
      $\sigma\sim_r\tau$ implies $t_a(\sigma)=t_a(\tau)$.
\end{enumerate}
The first condition implies the second by
\eqref{obs:free-behavior-stage}.  Conversely, each
$S_a(\rho_r)$ is surjective, since a finite stack of elements of
$P_r$ can be lifted through the surjection $\rho_r$.  For
$\eta\in S_a(P_r)$ define $u_a(\eta)$ to be $t_a(\sigma)$ for
any lift $\sigma$ of $\eta$.  Such a lift exists, and any two lifts
are $\sim_r$-equivalent by
\eqref{obs:free-stack-equivalence}.  The second condition therefore
makes the definition unambiguous.  It defines a unique
$u\in\operatorname{Tab}(P_r)=C_{r+1}$ with $\beta_r(u)=t$.
This proves the equivalence and uniqueness of the factor.

Consequently
\begin{equation}
  \beta(C_*)
    =\bigcup_{r\geq0}\beta_r(C_{r+1})
    =\left\{
       t\in\operatorname{Tab}(P_*):
       \begin{array}{l}
       \text{for some }r,\text{ each }t_a\text{ is constant}\\
       \text{on every }\sim_r\text{-class},\quad a\in\{c,e\}
       \end{array}
       \right\}.
  \label{obs:free-uniform-factorization}
\end{equation}
The one stage $r$ here is uniform over all stack lengths, keys,
payloads, and both kinds of ending.  The preceding argument proves
factorization for precisely this uniform condition.

\paragraph{Continue-prefix slicing.}
If $\sigma\in S_c(Y)$ and $t\in\operatorname{Tab}(Y)$, define
the slice $t|\sigma\in\operatorname{Tab}(Y)$ by
\begin{equation}
  (t|\sigma)_a(\tau)=t_a(\sigma\tau),
       \qquad \tau\in S_a(Y),\quad a\in\{c,e\},
  \label{obs:free-slice}
\end{equation}
where juxtaposition denotes concatenation.  The prefix has only
continue frames, so concatenation preserves the kind of ending.
This immediately gives slices of plus tables, by taking $Y=C_*$.

For minus tables, slicing preserves the representing stage.
For $u\in C_{r+1}$ and $\sigma\in S_c(P_*)$,
stack relabeling commutes with concatenation and gives
\begin{equation}
  \beta_r(u)|\sigma
     =\beta_r\bigl(u|S_c(\rho_r)(\sigma)\bigr).
  \label{obs:free-slice-bound}
\end{equation}
To check this equality, evaluate either side on
$\tau\in S_a(P_*)$; both values are
\[
  u_a\bigl(S_c(\rho_r)(\sigma)\,
                 S_a(\rho_r)(\tau)\bigr).
\]
The sliced table inside $\beta_r$ belongs to
$\operatorname{Tab}(P_r)=C_{r+1}$.  Thus a minus value represented
at $C_{r+1}$ has every continue-prefix slice represented at that
same stage.  Because $\beta$ is injective, this defines an intrinsic
slice on $C_*$, independent of the chosen representing stage.
Equivalently, fixing a prefix preserves invariance under
$\sim_r$ on all remaining payloads.  Default tables slice to
default tables.  Define
\[
 \mathcal K_T=\begin{cases}P_*,&T\text{ is free plus},\\
 C_*,&T\text{ is free minus}.\end{cases}
\]
Evaluate a minus value through the injective behavior map $\beta$.
Thus a plus value is an arbitrary table, whereas a proposed minus
behavior defines a value only after one uniform stage bound has been
proved.  The evaluation of a lambda at a free minus type will be the
place where this distinction matters.


\subsection{Finite typing trees and environments}
\label{obs:trees}

Evaluation must remember the particular product used to type an
application.  We therefore equip the finite syntax tree of a judgment
$\Gamma\vdash M:T$ with generation data.  At each node we store its
expected-type judgment and a core typing converting to that expectation:
the declaration at a variable; the axiom target at a sort; the two child
sorts and the formation triple at a product; the annotation sort, body
type and supporting product typing at a lambda; and the function's
exact product expectation and the argument's domain expectation at an
application.  We include the annotation's typing at a lambda.  These
data exist by generation; their witnesses are finite.  They do not
require uniqueness of sorts or formation rules.

The following operations on these trees will be used.  Valid thinning
and context conversion transport all judgments and generation witnesses.
At a variable the implicit core declaration is changed accordingly.
Retargeting a root changes its expected type to a convertible type at
which the same expression is typed; it retains the core data and all
children.  Capture-avoiding substitution transports the stored judgments
and grafts copies of the argument tree at occurrences of the substituted
variable, thinning and retargeting those copies as needed.  These
operations use only the previously established metatheory.

An environment for a finite generic context $\Gamma$ is a typed
substitution $\xi:\Gamma\to\Omega$, with arbitrary raw term images,
together with a parameter $\theta_x\in\mathcal K_{(A\xi)^\#}$ for each
declaration $x:A$ whose actual type $(A\xi)^\#$ is active.  Declaration
types are evaluated at their prefixes.  The notation $A\xi$ may also
use the whole substitution, since no later variable is free in $A$.
An extension by $x:=n$ is typed whenever $n$ has the actual normalized
domain type; conversion gives its required raw type.  At an active
domain it may be accompanied by any parameter in the corresponding
space.  At an inactive domain no parameter is supplied.

Two environments are \emph{compatible} if their raw images are
pointwise convertible and their corresponding parameters are equal.
Their actual normalized types are equal, so this equality of parameters
is well-typed.  All prefix changes are understood by thinning.

\subsection{Evaluation and reading candidates}
\label{obs:evaluation}

For a typing tree $t$ of $M:T$ and an environment $(\xi,\theta)$, define
\[
 \operatorname{Val}(t;\xi,\theta)\in\mathcal K_{(T\xi)^\#}
 \quad\text{when $(T\xi)^\#$ is active.}
\]
For a tree of $M:s$ with literal expected sort, define
$\operatorname{Read}(t;\xi,\theta)$ to be a candidate at the actual
expression $(M\xi)^\#$.  The two constructions are distinct: the
space index of the first is the actual \emph{type}; the candidate index
of the second is the actual \emph{expression}.

Every active space has a base in $\mathcal B$: its prefix base when
plain, and its empty-stack output when free.  Write
$\operatorname{base}(v)[A]$ for the candidate at actual type $A$ selected
by that base.

First specify the product candidate used in both constructions.
For a product tree with children $D:a,E:b$, put
\[
 D_0=(D\xi)^\#,
 \qquad A=((\Pi x:D.E)\xi)^\#,
 \qquad A_n=(E\xi[x:=n])^\#.
\]
The binder is chosen fresh for the images.
When $\pi(A)\in S$, define $\mathcal P(t;\xi,\theta)$ to consist of
the strongly normalizing $h\in\operatorname{Ty}(A)$ satisfying the
following test.  The tester $n$ ranges over
\[
 \begin{cases}
 \operatorname{Read}(D:a;\xi,\theta),&\pi(D_0)\in S,\\
 \operatorname{Ty}(D_0)\cap\mathrm{SN}_\beta,&\pi(D_0)\notin S.
 \end{cases}
\]
If $\pi(A_n)\in S$, require
\begin{equation}\label{obs:product-test}
 h\,n\in\operatorname{Read}(E:b;\xi[x:=n],\theta[x:=v])
\end{equation}
for every $v\in\mathcal K_{D_0}$ when $D_0$ is active, and for the
unique parameter-free extension otherwise.  If $\pi(A_n)\notin S$,
require just that $h\,n$ be strongly normalizing at its exact type
$A_n$.  Normalization and conversion identify $A_n$ with the
instantiation of the normalized product $A$.  The candidate test lemma
therefore shows that $\mathcal P$ is a candidate whenever its child
Reads are defined.

The clauses for $\operatorname{Val}$ are as follows.
\begin{enumerate}
\item A variable looks up its parameter.  A sort has the distinguished
      value of its actual type.
\item At a product with literal core sort $s$, start with the
      distinguished value.  If $\{s\}\in S$ and the profile of the
      actual product expression is in $S$, replace its base coordinate
      at that expression by $\mathcal P$.  All other coordinates and
      later observations stay distinguished.
\item At a lambda $\lambda x:D.P$, the base is $\mathbf t$.
      If its actual type is plain, prescribe its next value at every
      legal frame $(q,v)$ to be
      \begin{equation}\label{obs:lambda-equation}
       \operatorname{Val}(P;\xi[x:=q],\theta[x:=v]).
      \end{equation}
      The absent-payload convention applies when the domain is
      inactive.  Product compatibility gives the asserted target type.

      At a free actual type, give the following behavior on raw stacks.
      For a nonempty stack, check that its first key $q$ is a normal
      term at the actual normalized domain and that its next actual
      type is active.  Check its payload: absent at an inactive domain,
      an opposite-sign free value at a free domain, or a plain value
      with exactly the domain's type tag at a plain domain.  Finally
      require a continue flag exactly when the next type is free, and
      an exit flag exactly when it is plain.  Lemma~\ref{obs:graph-lemma}
      gives the same sign at a free next type and the opposite sign at
      a free domain.  On a valid continue frame, observe
      \eqref{obs:lambda-equation} on the remaining stack.  On a valid
      exit frame, return that value tagged with its plain target type.
      On an invalid stack, return the distinguished output for its
      ending.  The finite-bound lemma below proves that this behavior
      represents a member of the minus space when the lambda is minus.
\item At an application $F\,N$ with active actual result type, return
      its distinguished value if the function's actual type is
      inactive.  Otherwise use the frame with key $(N\xi)^\#$ and
      payload $\operatorname{Val}(N)$ when its actual domain is active.
      At a plain function take its next value.  At a free function take
      a continue slice when the result is free.  When the result is
      plain, take the exit output and project its correctly tagged
      value, using the distinguished target value for a missing or
      wrong tag.  The graph lemma and subject reduction make this a
      legal frame with the asserted target.
\end{enumerate}

The clauses for $\operatorname{Read}(M:s)$ are these.  If the actual
expression profile is outside $S$, or if $M$ is a literal sort, use the
top candidate.  At product syntax with actual profile in $S$, use
$\mathcal P$.  At all other syntax with actual profile in $S$, use
\begin{equation}\label{obs:read-base}
 \operatorname{base}(\operatorname{Val}(M:s))[\,(M\xi)^\#\,]
\end{equation}
if $\{s\}\in S$ and $\operatorname{Ax}(s)$ is active; use top otherwise.
The layer required for the Val in \eqref{obs:read-base} is exactly
$\operatorname{Ax}(s)$.

Val calls Read only on strict syntax children.  Read may call Val at
the same root.  Thus the intended recursion order is syntax size,
with Val preceding Read at a fixed size.  The next lemma establishes
totality as well as the only representability issue, at a minus lambda.
It also establishes compatible-image invariance without assuming that
the raw images are strongly normalizing.

\subsection{Uniform finite bounds for typed evaluation}
\label{obs:bound-section}

We verify that the preceding clauses define actual values in the stated
spaces, including at free minus layers.  In this subsection a typing tree
is fixed, but its typed substitution and its parameters may vary.  Bounds
will depend only on the finite tree and an integer, and not on any of
these images or parameters.  We use the images of $C_k$ in $C_*$ without
further notation.

Two environments for the same tree are called \emph{image-compatible} if their
term images are pointwise convertible.  Their actual normalized
declaration types consequently agree.  For $r\geq 0$, say that they
\emph{agree to level $r$} if, in addition, corresponding plain and minus
parameters are equal and corresponding plus parameters are
$\equiv_r$-equivalent.  Parameters at inactive declarations are absent.
An environment is \emph{minus-bounded by $k$} if all its minus parameters
belong to $C_k$.  In particular, being minus-bounded by zero means having
no minus parameters.  Every finite environment is minus-bounded by some
$k$, since $C_*$ is the increasing union of its stages.

\paragraph{The three binder facts.}
\label{obs:bound-binder-facts}
The following consequences of the graph definitions will be used
repeatedly.  Suppose a normalized product has profile $K$, domain profile
$I$, open codomain profile $J$, and instantiated codomain profile $L$.
If $K$ and $L$ are active, the positive ascent $L\longrightarrow
J\longrightarrow K$ makes $J$ active.  If also $I\in H$, the retained
negative edge $I\longrightarrow K$ makes $I$ active.  Therefore:
\begin{enumerate}
\item If $K$ is plain, every active frame payload and every active
continuation is plain: an active edge into a plain vertex starts at a
plain vertex, and the same holds for a positive path into a plain vertex.
\item If $K$ and $L$ are free, they have the same sign; a free domain
payload has the opposite sign to $K$.
\item If $K,L\in S$, then $J\in S$, since $J$ lies on a path between two
vertices of the primary strongly connected component $S$.  Hence an
active $I$ is direct.  A free direct vertex is plus, so such a domain is
never free minus.
\end{enumerate}
The last assertion is the reason that universal quantification over
payloads in a product candidate preserves a fixed minus bound.

\begin{lemma}[Uniform finite bounds]\label{obs:bound-statement}
For every fixed finite tree $t$ there are bounds for each of its
$\operatorname{Val}$ and $\operatorname{Read}$ evaluations, denoted below
by $F_t(k)$ when the evaluation is understood, with
$F_t(k)\geq k+1$, satisfying the following assertions.
\begin{enumerate}
\item Every evaluation is defined for every typed environment.
\item In any environment minus-bounded by $k$, a free minus
$\operatorname{Val}$ output belongs to $C_{F_t(k)}$.
\item For two environments minus-bounded by $k$ and agreeing to level
$F_t(k)$, plain and minus $\operatorname{Val}$ outputs are equal,
$\operatorname{Read}$ outputs are equal, and plus
$\operatorname{Val}$ outputs are $\equiv_k$-equivalent.
\end{enumerate}
The output spaces and their indices in the two environments agree:
pointwise convertibility of the images gives the same normalized actual
types, expression profiles, normalized argument keys, and plain type
tags.  Increasing a chosen bound preserves all three assertions.
\end{lemma}
\begin{proof}
We prove the statement together with totality by induction on syntax
size, proving $\operatorname{Val}$ before $\operatorname{Read}$ at a
given size.  All calls to $\operatorname{Read}$ from a product
$\operatorname{Val}$ concern proper children; the only same-size call
is from the remaining $\operatorname{Read}$ clause to
$\operatorname{Val}$.  Thus this is a well-founded order.  A bound for a
node is the maximum of the finitely many bounds described below for its
possible clauses.  Conditional variation of the actual profiles does
not introduce infinitely many cases: the estimates distinguish only
inactive, plain, plus, and minus layers and the finitely many syntactic
evaluation clauses.

\paragraph{Lookup and defaults.}
\label{obs:bound-leaves}
A minus lookup already lies in $C_k$, and a plus lookup agrees under
$\equiv_k$ whenever the environments agree to any level at least $k$.
Plain lookups agree exactly.  Constant-by-ending defaults lie in $C_1$
on the minus side.  Sort defaults and skipped-application defaults
therefore satisfy the assertions with bound $k+1$.

\paragraph{Lambdas at plain and plus layers.}
\label{obs:bound-positive-lambdas}
Write $l$ for the bound of the body evaluation at $k$.  A lambda at a
plain layer is defined using the plain prefix bijection.  Each of its
frames introduces an identical plain parameter, or no parameter, when
two matching frames are compared.  Its active continuation is also
plain, by the first binder fact.  The extended environments remain
minus-bounded by $k$, so body comparison gives equal frame outputs.
Together with the equal default base, the prefix bijection gives equal
lambda values.  This holds simultaneously for all frames and all
continuations; the same $l$ works for every one.

A lambda at a plus layer is a full table.  It is defined on every raw
stack using the gate in the evaluation definition.  For each individual
valid first frame, its minus payload, if present, belongs to some finite
stage of $C_*$.  Together with the finitely many original environment
parameters this gives a finite bound for the extended environment.
Totality of the smaller body tree therefore defines the requested body
value.  This use of totality does not require a single bound for all
possible payloads in the full table.

For comparison under $\equiv_k$, however, only stacks whose minus
payloads all lie in $C_k$ are inspected.  Compare the same such stack in
the two tables.  The gate agrees in the two environments.  If it is
valid, its first frame introduces either the same minus parameter in
$C_k$, the same plain parameter, or none.  Body comparison at $k$ thus
applies.  For a continue frame the body has plus sign and its
$\equiv_k$-agreement gives the same observation on the restricted tail.
For an exit frame the body is plain and agrees exactly.  Invalid stacks
and the empty stack have matching defaults.  This proves
$\equiv_k$-agreement of the plus lambda using the bound $l$.

\paragraph{Lambdas at minus layers.}
\label{obs:bound-minus-lambdas}
Again let $l$ be the body's bound at $k$, so $l\geq k+1\geq 1$.
First define the lambda's behavior on raw stacks by the stated gate and
body clauses.  Every valid first frame introduces a plus parameter, a
plain parameter, or none; no new minus parameter is introduced.  The
body is therefore total in an environment still minus-bounded by $k$.

Consider two original environments agreeing to level $l$, and two
stacks with identical keys, flags, none/plain payloads, and pairwise
$\equiv_l$-equivalent plus payloads.  Their gates agree.  At a valid
first frame the extended environments still agree to level $l$ and
remain minus-bounded by $k$.  If the frame continues, the two body
outputs are the same minus value $c\in C_l$.  Its behavior factors
through $\equiv_{l-1}$, and hence also through $\equiv_l$, on the plus
payloads of the tails.  It therefore gives equal observations on the
two tails.  If the frame exits, body comparison gives equal plain
values.  Defaults and empty observations also agree.

Apply this argument first with the same original environment on both
sides.  It proves uniform invariance of the lambda behavior under
$\equiv_l$ in all plus payloads.  The exact factorization result for
free spaces then represents it by an element of $C_{l+1}$.  Apply the
argument next to different original environments and the same stack on
both sides.  Their represented behaviors agree on every stack; the
behavior representation is injective, so their minus values are equal.
Thus $l+1$ is a valid output and comparison bound for the minus lambda.

\paragraph{Applications.}
\label{obs:bound-applications}
Let $f$ and $n$ be bounds for the function and argument evaluations,
respectively.  Actual normalized argument keys agree in image-compatible
environments.  A function at an inactive layer returns the default.  A
function at a plain layer has only plain active dependencies, by the
first binder fact.  The function and any evaluated argument payload
therefore agree exactly by their induction hypotheses at $k$, and the
plain prefix operation gives equal outputs.

Suppose the function is plus.  Argument comparison at $k$ makes a minus
payload exactly equal in the two environments and puts it in
$C_{n(k)}$.  Set
\[
                 a=\max\{k,n(k)\}.
\]
The original environments are also minus-bounded by $a$.  Function
comparison at $a$ gives $\equiv_a$-agreement, provided the environments
agree to level $f(a)$.  Every prefixed stack needed to compare a
continuing output under $\equiv_k$ has its new minus payload in
$C_{n(k)}\subseteq C_a$ and all tail minus payloads in
$C_k\subseteq C_a$.  The function observations therefore agree on
these stacks.  An exit is a single such frame and gives equal tagged
plain values, hence equal results after projection.  If the payload is
none or plain, the same estimate applies with no use of $C_{n(k)}$.

Suppose instead the function is minus.  Function comparison at $k$
gives the same value in $C_j$, where
\[
                         j=f(k).
\]
Apply argument comparison at $b=\max\{k,j\}$.  If the argument payload
is plus, the two payloads agree under $\equiv_b$, hence under
$\equiv_j$.  Since the function behavior already factors through
$\equiv_{j-1}$, inserting these payloads gives the same observations
on identical arbitrary tails.  A continue slice remains in $C_j$ by
the slicing result, and the two slices are equal by injectivity of the
behavior representation.  An exit gives equal plain values.  Equal
plain payloads and absent payloads are immediate special cases.

All these calls, together with the required output stages, are dominated
by the finite expression
\begin{equation}
  1+\max\bigl\{k,f(k),n(k),
       f(\max\{k,n(k)\}),n(\max\{k,f(k)\})\bigr\}.
  \label{obs:bound-application-formula}
\end{equation}
No monotonicity of $f$ or $n$ is needed: the displayed arguments are
the precise ones used above.  A plus active function can only have a
free plus active continuation, and a minus one only a free minus
active continuation, so the cases prove the required output assertion
as well as comparison.

\paragraph{Product candidates and product overrides.}
\label{obs:bound-products}
Consider a product candidate whose actual expression profile is in $S$.
Use the domain child's $\operatorname{Read}$ comparison at $k$ whenever
that Read is requested.  It gives exactly the same domain candidate in
the two environments, and hence exactly the same set of testers.  If
the domain Read is not requested, the typed strongly normalizing tester
set is already identical because the actual normalized domains agree.

For each common tester $n$, the instantiated actual codomain types
agree.  When their profile lies in $S$, the third binder fact says
that an active domain payload is plain or plus, never minus.  Pair
each payload choice with the identical choice in the other environment.
This preserves the minus bound $k$ and agreement at every level
already required for the original environments.  The codomain child's
$\operatorname{Read}$ comparison at $k$ therefore gives equal codomain
candidates for every tester and every payload choice.  Its bound is
uniform in the raw term $n$ and the payload: both are variable images
or parameters for the same fixed smaller tree.  If the instantiated
profile is outside $S$, the required typed strong-normalization
condition is identical on the two sides.  Thus all the universally
quantified membership conditions defining the product candidate agree,
so the candidate itself is equal.  The maximum of the two child bounds
and $k+1$ suffices.

The product $\operatorname{Val}$ override places this candidate in one
component of the empty-stack $\mathcal B$ observation, leaving the
other components and all nonempty observations at their defaults.
The actual component index is the same in image-compatible environments.
For a minus output such a table is represented already in $C_1$,
because its nonempty observations are constant by ending and its empty
observation contains no payload.  For a plus output its observations
agree exactly, in particular under $\equiv_k$.  At a plain layer the
default continuations and equal base give equality by the prefix
bijection.  This proves the override case with the same child bounds.

\paragraph{The remaining Read clauses.}
\label{obs:bound-read}
The default Read clauses are immediate.  At a nonproduct whose Read
uses its $\operatorname{Val}$ base, use the already proved
same-size $\operatorname{Val}$ bound.  Plain and minus values agree
exactly; plus values agree under $\equiv_k$.  The latter also gives
equal empty observations, even when $k=0$, since the empty stack
belongs to every restriction.  Reading the same actual-expression
component therefore gives the same candidate.  This completes the
simultaneous induction.
\end{proof}

\paragraph{Exact invariance for identical parameters.}
\label{obs:bound-exact-invariance}
As a consequence, image-compatible environments with identical parameters
give identical outputs of every kind.  For plain, minus, and Read
outputs choose a common finite minus bound and use the comparison
statement.  For plus outputs, fix any raw observation stack.  Its
finitely many minus payloads and the finitely many minus parameters of
the original environments lie in some common $C_k$.  The environments
agree to every level because their parameters are identical, so the
comparison statement gives $\equiv_k$-agreement of the plus outputs
and hence equality on this stack.  The stack was arbitrary, proving
equality of the full tables.

The finite-bound proof applies to arbitrary typed trees; it assumes
neither normal syntax nor goodness.  In particular its uniformity is
unaffected by the lengths of reductions needed to normalize varying
raw images.  Nonfunctional axioms and product rules also cause no new
case: the fixed tree stores its selected judgments, while compatibility
aligns the corresponding actual normalized indices.


\begin{lemma}[Structural invariance]\label{obs:structural-invariance}
Compatible environments give identical Val and Read outputs.  Valid
thinning and context conversion, with transported trees and compatible
old environments, preserve both outputs.  Retargeting a root to a
convertible expected type preserves Val.  All these operations commute with capture-avoiding renaming of generic
declared and bound variables, reindexing the substitution and parameters
while leaving their actual images and values unchanged.
\end{lemma}
\begin{proof}
For compatible environments, apply the finite-bound lemma.  Plain,
minus, and Read equality follows at once.  For plus equality, fix any
finite raw stack.  Its finitely many minus payloads lie in some common
$C_k$; equality through level $k$ gives equality on that stack.  Thus
the tables are equal.

For the structural operations, induct in the Val/Read syntax order.
Actual normalized types and frame keys agree.  At a binder, the raw
domain types are equal or convertible sorted types, so the same actual
term gives a typed extension on both sides and the induction hypothesis
applies for every matching payload.  This proves equality of all
prefix coordinates or raw observations.  At a product candidate it
proves equality of domain tests and, for every raw tester and payload,
of the codomain candidates.  The application clause then extracts the
same coordinate or slice.  The other clauses are immediate from lookup,
defaults, and base projection.  Retargeting changes no core clause or
child and leaves the actual index unchanged.  Bound variables may be chosen fresh throughout. Reindexing a generic
declaration leaves its actual image, parameter, and every subsequent actual
type unchanged; the same induction therefore also proves equivariance
under renaming generic declarations.
\end{proof}

\subsection{Good evaluations and substitution}
\label{obs:goodness}

Arbitrary stored trees need not respect substitution: the default used
at an inactive function can hide syntax that later becomes relevant.
We identify the trees for which this cannot occur.  Goodness will
depend on the raw images $\xi$, but never on their parameters $\theta$.

For an active Val evaluation, define goodness by the following clauses.
A lambda is good if its body is good under every typed raw extension
$x:=n$ for which the actual next type is active.  An application with
inactive actual function layer is good if its syntactic head is a
variable.  An application with active function layer is good if the
function is good and the argument is good whenever its actual layer is
active.  A product whose base is overridden is good if its product
candidate is good.  Lookups and all other default clauses are good
without further conditions.

The goodness of a used product candidate means that its domain Read is
good when used, and its codomain Read is good under \emph{every} typed
raw extension having actual codomain profile in $S$.  There is no
candidate-membership restriction on these extensions.  For Read,
defaults outside actual profile $S$, and sorts, are good.  A product
Read in $S$ is good when its product candidate is good.  A remaining
Read in $S$, with expected sort $s$, is good if $\{s\}\in S$ and,
when $\operatorname{Ax}(s)$ is inactive, its syntactic head is a
variable; when that layer is active, its Val must be good.
These are again definitions by the Val/Read syntax order, and their
binder clauses quantify over raw terms, not over reduction sequences.

\begin{lemma}[Basic goodness facts]\label{obs:goodness-basic}
Goodness has the structural invariances of
Lemma~\ref{obs:structural-invariance}.  If the actual expression profile
and $\{s\}$ are in $S$ and $\operatorname{Ax}(s)$ is active, then
Read$(M:s)$ is always the base projection
\eqref{obs:read-base}, at every syntactic shape, and its goodness is
equivalent to Val goodness.

Every normal syntax tree is good for active Val, under every typed raw
substitution.  A normal tree $M:s$ is good for Read whenever its generic
expression profile belongs to $U\cup S$.
\end{lemma}
\begin{proof}
Structural invariance follows by the same recursion as for evaluation,
matching all typed raw binder extensions by conversion.  For the base
projection assertion, the sort case gives top on both sides, the
product case is exactly the override, and all other cases use
\eqref{obs:read-base}.  In the product case its literal core sort equals
$s$ because it is convertible to $s$.  The goodness clauses agree for
the same reason.

Prove the normal-syntax claims simultaneously in the Val/Read syntax
order and universally over typed raw substitutions.  Normal applications
are variable-headed, so skipped ones are good.  Lambda bodies are
normal and the induction hypothesis applies to every typed extension.
At an overridden normal product, its generic profile contains the
core sort $s$, with $\{s\}\in S$.  The positive inclusion edge places
that generic profile, and hence both child profiles, in $U\cup S$.
The child Read induction hypotheses prove all required goodness
conditions.  The same predecessor argument applies at product Read.
For a nonsort, nonproduct normal Read whose actual profile is in $S$,
the generic profile cannot be in $U$, since growth from $U$ stays in
$U$.  It is therefore in $S$.  The normal expression is neutral and
has singleton generic profile $\{s\}$, giving exactly the required
Read condition; its Val, when needed, is good by induction.
\end{proof}

\begin{lemma}[Substitution of a good value]\label{obs:substitution}
Let a tree over $\Gamma,x:D,\Delta$ be given, and let $u$ be a tree of
$\Gamma\vdash N:D$.  Substitute $u$ capture-avoidably at occurrences
of $x$, transporting the context, expectations, and generation data.
Give the resulting context a typed environment, and give the original
context a compatible environment whose image of $x$ is convertible to
the actual expression $N\xi$.  Require the remaining corresponding
parameters to agree.

Let $I=\pi((D\xi)^\#)$.  Assume that $I$ is active whenever $I\in H$.
If $I$ is active, assume Val$(u;\xi)$ is good and use its value as the
original parameter of $x$.  Then a good Val or Read evaluation before
substitution is good after substitution, and the two outputs are equal.
\end{lemma}
\begin{proof}
Composition of capture-avoiding substitutions makes corresponding
actual images convertible and corresponding actual types equal.
Induct on the original tree in the Val/Read order.  An active lookup
of $x$ is exactly the chosen value of $u$; thinning and retargeting
preserve it by structural invariance.  Other lookups agree.

At a binder, rename its variable fresh for $N$.  The same raw binder
image and payload extend both environments, the hypotheses on $I$
remain unchanged, and the induction hypothesis applies.  This gives
equality at every legal frame or every raw observation and proves all
universally quantified goodness requirements.  At a product the root
persists.  Domain Reads agree by induction; for every raw tester and
payload the codomain Reads agree by the binder argument.  Thus the
entire product candidates agree and remain good.  At an unskipped
application, the evaluated children agree by induction and its normal
argument keys agree by conversion, so it extracts the same output.

It remains to justify the skipped clauses, where an arbitrary change of
head would invalidate the argument.  Suppose a skipped, variable-headed
application had head $x$.  Trace its successive actual typing products
backward along the spine.  There is a primary positive path from its
active result layer to $I$, passing through its supposedly inactive
function layer.  The path stays in $H$ by positive closure.  By
hypothesis $I$ is active; consequently every vertex on the path reaches
a direct vertex and is active.  This is a contradiction.  The head is
therefore a different variable and persists under substitution, so the
skipped application stays good and keeps its default.

For a skipped nonproduct Read, goodness supplies $\{s\}\in S$ and a
variable head.  If that head were $x$, the same spine argument would
start at $\operatorname{Ax}(s)$.  This is nonempty because a neutral
typed at $s$ has sorted $s$, and it is a seed in $H$.  It would reach
the active vertex $I$, forcing $\operatorname{Ax}(s)$ active, contrary
to this being a skipped Read.  Thus that head also persists.  In the
active case use the all-shape base projection of
Lemma~\ref{obs:goodness-basic} and the Val induction already proved,
even if substitution changes the root shape.  Read outside actual
profile $S$ stays the same default.  This exhausts the clauses.
\end{proof}

\subsection{Reduction and normal-tree independence}
\label{obs:reduction}

\begin{lemma}[Reduction transport]\label{obs:reduction-transport}
A one-step full beta reduction of a typed expression transports its
tree to a tree of the reduct at the same expected type.  For every
environment, good Val and Read evaluations remain good and their
outputs are preserved.  This includes reductions in annotations.
\end{lemma}
\begin{proof}
Construct the transported tree independently of its environment.  At
contextual steps transport the changed child.  A change in a product
domain context-converts its open codomain tree.  A change in a lambda
annotation also transports the supporting product typing by subject
reduction and context-converts the open body.  A change in an
application argument changes its core result type to a convertible
one; subject reduction retains the root expectation.

Induct on the reduction position, simultaneously for Val and Read.
At evaluated children use the induction hypothesis.  At binders apply
it under every typed raw extension required by goodness; structural
invariance handles the context-converted unchanged children.  Actual
types, gates, and argument keys are unchanged.  At a skipped
variable-headed spine, a reduction can only be inside its arguments;
the head and skipped condition persist.  Ignored children need no
semantic property, only the transported typing.  For a remaining Read
with active base projection, apply the Val statement already proved
at the same root.  A Read outside actual profile $S$ remains default.

At a root contraction $(\lambda x:D.P)N\longrightarrow P[N/x]$,
generation and product compatibility align the lambda annotation and
body type with the application's stored product.  Retarget the
argument to the annotation type, substitute its tree in the body, and
retarget the result to the old expectation.

Consider a good active Val.  The application cannot have been skipped,
since its syntactic head is a lambda.  Let $K,L$ be the actual function
and result layers, respectively.  They are active.  By
Lemma~\ref{obs:graph-lemma}, the actual domain layer $I$ is active
whenever it belongs to $H$.  The argument is good if $I$ is active.
Evaluation uses the valid normal key $(N\xi)^\#$ and, when required,
the payload Val$(N)$.  The lambda equation gives its body value with
this key, including at an exit where the lambda supplied the correct
tag.  The raw image $N\xi$ is convertible to its normal form.
Structural invariance therefore identifies this body value with the
one at the raw image, with the same payload.  Lambda goodness supplies
body goodness for both typed extensions.  Lemma~\ref{obs:substitution}
now identifies that value with the beta contractum and proves its
goodness.  No strong normalization of $N\xi$ was used.

For a Read with actual profile in $S$, a skipped clause at the root is
impossible because the head is a lambda.  Its goodness therefore
provides the active all-shape base projection, and the preceding Val
argument applies.  This proves the root case and the lemma.
\end{proof}

\begin{lemma}[Independence on normal syntax]\label{obs:normal-independence}
On a fixed normal expression in a fixed generic context, Val is
independent of its stored tree whenever the root expectations are
convertible.  Read is independent of the stored tree and of its chosen
literal expected sort.
\end{lemma}
\begin{proof}
Induct in the Val/Read syntax order, comparing the same environments.
Variables use the same declaration parameter and sorts use defaults.
For product Val, convertible expected types force its two literal core
sorts to agree.  The override condition is therefore identical.  For
product Read, the product candidate compares the child Reads
independently of their sort choices, by induction.  Its raw extensions
and actual indices agree, so the candidates are equal.

For a lambda, product compatibility aligns the body expectations up to
conversion.  The annotation is fixed syntax.  The body Val induction
therefore compares every frame and observation.  In a normal
application the function is variable-headed, so its generic type is
unique up to conversion.  Product compatibility aligns the domains and
codomains in the two application trees.  The function and argument
Val induction hypotheses give the same keys, payloads, and selected
outputs.  Finally, a normal nonsort, nonproduct expression is neutral
and has a unique literal sort.  Thus any nondefault Read uses the same
Val base.  These are all normal syntactic shapes.  Only neutral type
uniqueness has been used; functionality of the specification is not
assumed.
\end{proof}

\subsection{The candidate interface}
\label{obs:interface}

We can now return from raw observations to the candidate interpretation.
For a generic normalized type $T$ with $\pi(T)\in S$, choose any
literal sort typing of $T$ and put
\begin{equation}\label{obs:R-definition}
        R_T(\xi,\theta)=\operatorname{Read}(T;\xi,\theta).
\end{equation}
It is independent of the choice by
Lemma~\ref{obs:normal-independence}.  At generic profiles outside $S$
use the typed-SN top, as in the candidate interface.  At actual
profiles in $U$, all candidates are top by the already established
equality with $G$.

\begin{proposition}[Observation interpretation]\label{obs:interface-theorem}
Assume system-wide weak normalization and \eqref{obs:forbidden}.
Then \eqref{obs:R-definition} gives candidates with the required
weakening and context-conversion invariance and nonempty parameter
extensions.  It has the product test clause and the specialization
property: if $T=\Pi x:D.E$ is generic normalized of profile in $S$,
$\Gamma\vdash N:D$, and $N\xi$ passes its domain test, then the test
of a candidate member $h$ entails
\[
                   h\,(N\xi)\in R_{(E[N/x])^\#}(\xi,\theta).
\]
The ordinary typed-SN clause applies when the instantiated actual
profile lies in $U$.
\end{proposition}
\begin{proof}
Candidate membership and the product clause are the Read construction
and the candidate test lemma.  The interpretation's invariances follow
from structural invariance and normal-tree independence.  Every active
space has its distinguished member, so arbitrary context extensions
can always be supplied with parameters.

For specialization, if the actual product profile is in $U$, its
candidate is the top and hence equals $G$.  The domain tester is
strongly normalizing.  Propagation of $G$ along this typed argument
gives the required typed strong normalization at the instantiated
profile, which also lies in $U$.  We may therefore suppose the actual
product profile $K$ is in $S$.  First suppose that the actual
instantiated codomain profile $L$ is in $S$, and let $I$ be the actual
domain profile.  Lemma~\ref{obs:graph-lemma} places the open codomain profile in $S$.
Thus if $I\in H$, the definition makes $I$ direct and hence active;
in particular, every active $I$ is direct.  Normalize the generic argument
$N$ to $N^\#$, retaining its type $D$ by subject reduction.  If $I$ is
active, choose as binder payload
\[
                     v=\operatorname{Val}(N^\#:D;\xi,\theta).
\]
This exists and is good by normal-syntax goodness.  The product test
quantifies over every payload and therefore includes this choice.

Its open codomain $E$ is normal, with generic profile in $U\cup S$,
so its Read is good.  Apply Lemma~\ref{obs:substitution} to substitute
$N^\#$ there, using the raw image $N\xi$.  This image is convertible
to $N^\#\xi$, as required.  The lemma identifies the tested open
codomain Read with the good Read of the generic expression
$E[N^\#/x]$.  This expression is typed by substitution and hence
weakly normalizing by the system hypothesis.  Transport its tree along
a normalizing sequence using Lemma~\ref{obs:reduction-transport}.
Its normal form is precisely $(E[N/x])^\#$, because $N$ and $N^\#$
are convertible.  Normal-tree independence identifies the final Read
with \eqref{obs:R-definition} at that type.

For completeness, its generic profile belongs to $U\cup S$: it has a
positive ascent to the original product profile.  It cannot belong
to $U$ when its actual profile is $L\in S$, because profile growth
from $U$ stays in $U$.  Thus \eqref{obs:R-definition} is indeed the
appropriate generic $S$ entry.  If the actual instantiated profile is
in $U$, the test demands typed strong normalization and the processed
component identity identifies this with the required top candidate.
These are the only actual profiles reached from the current component.
\end{proof}

Nothing in this construction requires $H$ to be nonempty.  If it is
empty there are no active spaces or parameters.  A normal neutral Read
of generic profile in $S$ would have profile $\{s\}$ with sorted $s$,
which would supply the nonempty seed $\operatorname{Ax}(s)$ in $H$.
Thus that case cannot occur, and the sort/product clauses complete the
same construction without an extra assumption.

The uses of weak normalization in this section are now explicit.  It
provides the normal forms of actual types, argument keys, and raw typed
substitutions, and it normalizes the generic expressions used in the
specialization proof.  Strong normalization enters only through the
candidate tests already specified in the normalization interface.  The
remaining task is to derive \eqref{obs:forbidden} from weak normalization.
